\section{Endereços, entropia, descrição e busca}
\label{sec:complexity}

\subsection{Comprimento do índice canônico}

Para um átomo com parâmetros $(k,m,y_0,s)$, defina
\begin{equation}
 H=\max\{2,\ k+1,\ |m|+1,\ \height(y_0),\ \height(s)\}.
 \label{eq:atom-height}
\end{equation}
Escreva $\Library_H\subset\Library$ para o conjunto dos átomos de altura no máximo $H$.

\begin{theorem}[cota polinomial do endereço]\label{thm:index-bound}
O único endereço $n$ do átomo satisfaz
\begin{equation}
 n=O(H^8),\qquad \log_2(n+1)=O(\log H).
 \label{eq:index-bound}
\end{equation}
As constantes são universais para a enumeração fixada.
\end{theorem}

\begin{proof}
Pelo \cref{lem:rational-rank}, os postos $r$ de $y_0$ e $t$ de $s$ são $O(H^2)$. Pelo \cref{lem:integers}, $\mu=Z^{-1}(m)=O(H)$, e $k=O(H)$. A cota do \cref{lem:cantor} dá
\[
 u=\pi(k,\mu)=O(H^2),\qquad v=\pi(r,t)=O(H^4),
\]
e, por uma última aplicação,
\[
 n=\pi(u,v)=O\bigl((u+v+1)^2\bigr)=O(H^8).
\]
Como $H\ge2$, tomar logaritmos fornece a segunda afirmação.
\end{proof}

\subsection{Contagem por altura e consequência minimax para o endereço}

O expoente $8$ do \cref{thm:index-bound} é um artefato do aninhamento triplo do pareamento de Cantor, não uma propriedade da biblioteca. A contagem a seguir separa as duas coisas.

\begin{lemma}[cardinalidade das cascas de altura]\label{lem:atom-count}
Para todo $H\ge2$,
\begin{equation}
 |\Library_H|\le18H^6.
 \label{eq:atom-count-upper}
\end{equation}
Além disso, existem uma constante universal $c_\Library>0$ e $H_1\ge2$ tais
que, para todo $H\ge H_1$,
\begin{equation}
 |\Library_H|\ge c_\Library H^6.
 \label{eq:atom-count}
\end{equation}
Em particular, $|\Library_H|=\Theta(H^6)$.
\end{lemma}

\begin{proof}
Pelo \cref{thm:dictionary}, um átomo é determinado bijetivamente pela quádrupla $(k,m,y_0,s)\in\Nzero\times\mathbb{Z}\times\mathbb{Q}\times\Qstar$, e a altura \eqref{eq:atom-height} limita cada coordenada separadamente. Assim
\[
 |\Library_H|=\#\{k\le H-1\}\cdot\#\{|m|\le H-1\}\cdot|\mathcal{Q}_H|\cdot|\mathcal{Q}_H\setminus\{0\}|,
\]
com $\mathcal{Q}_H$ como no \cref{lem:rational-count}. O primeiro fator é $H$, o segundo é $2H-1$, e os dois últimos estão entre $c_1H^2$ e $3H^2$ para $H$ grande. Logo
\[
 |\Library_H|\le H\cdot 2H\cdot 3H^2\cdot 3H^2=18H^6,
\]
\[
 |\Library_H|\ge H\cdot H\cdot c_1H^2\cdot\bigl(c_1H^2-1\bigr)\ \ge\ \frac{c_1^2}{2}H^6
\]
para todo $H$ suficientemente grande. Tome $c_\Library=c_1^2/2$.
\end{proof}

\paragraph{Dependência do expoente no calibre.}
O expoente $6$ pertence à exaustão definida por \eqref{eq:atom-height}, não à
biblioteca sem calibre. De fato, se os quatro orçamentos forem desacoplados,
\[
 \Library_{K_0,M_0,Y,S}
 :=\{(k,m,y_0,s):k\le K_0,\ |m|\le M_0,
       \ \height(y_0)\le Y,\ \height(s)\le S\},
\]
então a bijeção do \cref{thm:dictionary} dá exatamente
\[
 |\Library_{K_0,M_0,Y,S}|
 =(K_0+1)(2M_0+1)|\mathcal Q_Y|
   (|\mathcal Q_S|-1).
\]
Assim, até constantes universais e para $Y,S$ grandes, a contagem é da ordem
de $(K_0+1)(M_0+1)Y^2S^2$. O calibre
\eqref{eq:atom-height} acopla $K_0=M_0=H-1$ e $Y=S=H$, produzindo $H^6$;
o acoplamento $K_0=M_0=\lfloor\log_2H\rfloor$ e $Y=S=H$, por exemplo,
produziria $H^4(\log H)^2$. Portanto mudar o orçamento de resolução, extensão
espacial ou precisão aritmética muda o expoente.

\begin{proposition}[ordem exata do pior endereço aritmético]
\label{prop:gamma-exact-order}
Se
\[
 M_\Gamma(H):=\max_{a\in\Library_H}\Gamma^{-1}(a),
\]
então
\begin{equation}
 M_\Gamma(H)=\Theta(H^8).
 \label{eq:gamma-exact-order}
\end{equation}
Consequentemente, o comprimento da maior palavra do gerador aritmético sobre
$\Library_H$ é $8\log_2H+O(1)$ bits.
\end{proposition}

\begin{proof}
A cota superior é o \cref{thm:index-bound}. Para a inferior, escreva
$N_H=|\mathcal Q_H|$. Pelo \cref{lem:rational-count}, $N_H\ge c_1H^2$ para
$H$ suficientemente grande. Como as enumerações de \eqref{eq:q-enums} listam
por altura, existem $y_H\in\mathcal Q_H$ e
$s_H\in\mathcal Q_H\setminus\{0\}$ cujos postos respectivos satisfazem
\[
 r_H=N_H-1,\qquad t_H=N_H-2,
\]
e, portanto, $r_H,t_H\ge c_2H^2$ para alguma constante universal $c_2>0$.
Considere o átomo de parâmetros $(k,m,y_0,s)=(0,0,y_H,s_H)$, que pertence a
$\Library_H$. Como
\[
 \pi(x,y)=\frac{(x+y)(x+y+1)}2+y\ge\frac{(x+y)^2}{2},
\]
temos $v_H=\pi(r_H,t_H)\ge c_3H^4$. Além disso,
$\pi(0,0)=0$, de modo que o endereço desse átomo satisfaz
\[
 n_H=\pi\bigl(\pi(0,0),v_H\bigr)=\pi(0,v_H)
 \ge\frac{v_H^2}{2}\ge c_4H^8.
\]
Logo $M_\Gamma(H)=\Omega(H^8)$, que, junto da cota superior, prova
\eqref{eq:gamma-exact-order}. A afirmação em bits segue tomando logaritmos.
\end{proof}

\begin{theorem}[cota minimax relativa ao calibre]\label{thm:address-lower}
Seja $A:\Library\to\Nzero$ uma aplicação injetiva qualquer. Para todo $H\ge2$,
\begin{equation}
 \max_{a\in\Library_H}A(a)\ge |\Library_H|-1.
 \label{eq:address-pigeonhole}
\end{equation}
Consequentemente, para todo $H\ge H_1$, existe um átomo de altura no máximo
$H$ com
\begin{equation}
 A(\text{átomo})\ \ge\ c_\Library H^6-1,
 \qquad\text{logo}\qquad
 \log_2\bigl(A(\text{átomo})+1\bigr)\ \ge\ 6\log_2H-O(1).
 \label{eq:address-lower}
\end{equation}
Em particular, nenhum esquema de endereçamento injetivo, canônico ou não,
pode usar $o(\log H)$ bits no pior caso sobre $\Library_H$. O expoente $6$ é
o expoente minimax para a exaustão determinada por \eqref{eq:atom-height}; não
é reivindicado como invariante sob outro calibre.
\end{theorem}

\begin{proof}
Por injetividade, $A(\Library_H)$ é um subconjunto de $\Nzero$ com
$|\Library_H|$ elementos distintos; seu máximo é, portanto, pelo menos
$|\Library_H|-1$, o que prova \eqref{eq:address-pigeonhole}. Aplique então
\eqref{eq:atom-count} e tome logaritmos.
\end{proof}

O \cref{thm:address-lower} é o princípio da casa dos pombos aplicado à
exaustão finita. O insumo matemático substantivo é a contagem do
\cref{lem:atom-count}; a realização exata abaixo decorre de listar uma
exaustão encaixada casca a casca. Essa separação é importante: a contribuição
é a contabilidade do objeto aritmetizado, não uma nova desigualdade de
codificação.

\begin{remark}
O argumento é puramente de contagem e não usa nenhuma propriedade de $\Gamma$. Ele vale igualmente para endereçamentos que dependam de tabelas, de estado ou de escolhas arbitrárias. Combinado com a \cref{prop:gamma-exact-order}, mostra que o gerador aritmético $\Gamma$ usa, no pior caso, $8\log_2H+O(1)$ bits, enquanto o limite minimax é $6\log_2H+O(1)$. Assim, o fator assintótico $8/6=4/3$ no comprimento da palavra está demonstrado, e não apenas sugerido por uma cota superior.
\end{remark}

Resta saber se o expoente $6$ desse calibre é atingível por uma enumeração
explícita. É, por uma construção que reutiliza a mesma ideia de casca.

\begin{definition}[enumeração por altura da biblioteca]\label{def:sharp}
Para $H\ge2$, a casca $\Library_H\setminus\Library_{H-1}$ é finita e não vazia. Ordene $\Library$ listando as cascas em ordem crescente de $H$ e, dentro de cada casca, ordenando as quádruplas $(k,m,y_0,s)$ lexicograficamente por $\bigl(k,\ Z^{-1}(m),\ \rank_Q(y_0),\ \rank_{Q^*}(s)\bigr)$, onde $\rank$ denota o posto nas enumerações \eqref{eq:q-enums}. A bijeção resultante é denotada $\Gamma^\sharp:\Nzero\to\Library$, e $n^\sharp$ é o endereço de um átomo nessa ordem.
\end{definition}

\begin{theorem}[endereço minimax exato para a exaustão]\label{thm:sharp-optimal}
A aplicação $\Gamma^\sharp$ é uma bijeção de $\Nzero$ sobre $\Library$, computável nos dois sentidos, e para todo $H\ge2$
\begin{equation}
 \max_{a\in\Library_H}n^\sharp(a)=|\Library_H|-1\ \le\ 18H^6-1 .
 \label{eq:sharp-upper}
\end{equation}
Além disso, para \emph{qualquer} endereçamento injetivo $A:\Library\to\Nzero$,
\begin{equation}
 \max_{a\in\Library_H}A(a)\ \ge\ |\Library_H|-1 .
 \label{eq:sharp-minimax}
\end{equation}
Logo $\Gamma^\sharp$ realiza, com igualdade e para todo $H$, o mínimo possível
do maior índice no pior caso sobre $\Library_H$; em particular esse pior
índice é $\Theta(H^6)$ e seu comprimento binário é minimax-ótimo para o
calibre \eqref{eq:atom-height}.
\end{theorem}

\begin{proof}
Cada casca é finita pelo \cref{lem:atom-count} e a ordem dentro dela é total, de modo que a concatenação está bem definida e é bijetiva; a computabilidade nos dois sentidos segue de as cascas serem finitas e explicitamente enumeráveis. Por construção, $\Gamma^\sharp$ lista primeiro, e exatamente, os elementos de $\Library_H$; portanto $\Gamma^\sharp$ restrita a $\{0,\ldots,|\Library_H|-1\}$ é uma bijeção sobre $\Library_H$, o que dá a igualdade em \eqref{eq:sharp-upper}, e \eqref{eq:atom-count-upper} dá a cota numérica. A desigualdade \eqref{eq:sharp-minimax} é \eqref{eq:address-pigeonhole}: a imagem $A(\Library_H)$ tem $|\Library_H|$ elementos distintos de $\Nzero$, logo seu máximo é ao menos $|\Library_H|-1$.
\end{proof}

\begin{remark}[qual endereço usar]\label{rem:which-address}
As duas enumerações resolvem problemas diferentes e ambas são mantidas neste trabalho.
\begin{itemize}[leftmargin=1.6em]
 \item $\Gamma$ tem inversa em \emph{forma fechada}: descompactar $n$ custa três raízes quadradas inteiras e um número constante de operações aritméticas, sem qualquer enumeração auxiliar. Seu preço é o expoente $8$.
 \item $\Gamma^\sharp$ atinge o mínimo exato de \eqref{eq:sharp-minimax}, logo o expoente minimax $6$ para o maior índice sob o calibre fixado, mas sua inversa requer calcular cardinalidades de cascas, isto é, somas de totientes truncadas; o custo é polinomial em $H$, não constante.
\end{itemize}
Em bits, os dois esquemas coincidem em ordem: $\log_2(n+1)=\Theta(\log H)$ em ambos, e o \cref{thm:address-lower} mostra que nenhum outro esquema faz melhor no pior caso sobre $\Library_H$. A diferença é uma constante multiplicativa $4/3$ no comprimento da palavra.
\end{remark}

Operacionalmente, $\Gamma$ é a semântica do átomo: seu índice é descompactado
diretamente em célula, valor ancorado e inclinação, e por isso é o endereço
usado em avaliações, máscaras e desempates canônicos. Já $\Gamma^\sharp$ é um

\emph{codec} para armazenamento ou transmissão sob um limite de altura
conhecido. Converter um endereço de $\Gamma$ para $\Gamma^\sharp$ significa
decodificar primeiro o átomo e recodificar a mesma quádrupla na ordem por
cascas; a conversão inversa faz o percurso oposto. Em ambos os sentidos, a
célula e a função afim selecionada permanecem exatamente as mesmas. Portanto,
a vantagem minimax de $\Gamma^\sharp$ não altera a semântica definida por
$\Gamma$.

Sob uma raiz, um domínio e uma convenção de fronteira compartilhados, os
endereços dos átomos já recuperam suas células pela
\cref{lem:cell-recovery}. O conjunto das células-folha determina a árvore
diádica podada ao fechá-lo por ancestrais até a raiz; a ordem espacial é
recuperada ordenando os extremos. Assim, basta transmitir $N$, $H$ e a lista
de endereços: um esqueleto separado seria redundante.

\begin{theorem}[serialização condicionada de montagens]\label{thm:assembly-serialization}
Fixe, como informação compartilhada, o domínio $K$, a raiz diádica inicial,
a convenção de fronteira e a enumeração usada. Seja $g\in\Aclass(K)$ descrita
por uma árvore diádica com $N$ folhas, e seja $H\ge2$ a maior altura de seus
átomos. Existe um fluxo de bits autodelimitado, decodificável sem conhecer
previamente $N$ ou $H$, com comprimento
\begin{align}
 L(g\mid K)
 &\le \bigl(2\lfloor\log_2N\rfloor+1\bigr)
      +\bigl(2\lfloor\log_2H\rfloor+1\bigr)
      +N\left\lceil\log_2(18H^6)\right\rceil \notag\\
 &=N\bigl(6\log_2H+O(1)\bigr)+O(\log N+\log H)
 \label{eq:assembly-bitrate}
\end{align}
quando os átomos são transmitidos por $\Gamma^\sharp$. Com o gerador
aritmético $\Gamma$, o termo por folha é $8\log_2H+O(1)$ bits.
\end{theorem}

\begin{proof}
Codifique $N$ e $H$, nessa ordem, pelo código gama de Elias, com comprimentos
$2\lfloor\log_2N\rfloor+1$ e
$2\lfloor\log_2H\rfloor+1$. Conhecido $H$, cada endereço de $\Gamma^\sharp$
ocupa exatamente
$b_H=\lceil\log_2(18H^6)\rceil$ bits, pois
$n^\sharp<18H^6$ pelo \cref{thm:sharp-optimal}. O cabeçalho $N$ informa
quantos blocos de largura $b_H$ devem ser lidos. Os endereços recuperam os
átomos e suas células; ordenar os extremos fornece a ordem espacial, e fechar
as células por ancestrais recupera a árvore podada. O domínio compartilhado
determina apenas o recorte nas extremidades. Isso prova a primeira cota e a
decodificação sem esqueleto separado. Para $\Gamma$, use
$\max_{a\in\Library_H}\Gamma^{-1}(a)=\Theta(H^8)$ da
\cref{prop:gamma-exact-order}.
\end{proof}

\begin{remark}
Sem o domínio e a convenção de fronteira compartilhados, a expressão
``descrição completa'' seria incorreta para um intervalo $K$ com extremos
reais arbitrários: esses dados têm custo próprio e não são codificáveis, em
geral, por uma palavra finita. O \cref{thm:assembly-serialization} explicita
precisamente a informação lateral sob a qual a taxa de bits é afirmada.
\end{remark}

\subsection{Três formatos prefix-free e propagação de altura}
\label{subsec:coding-formats}

As comparações de bits a seguir usam fluxos concretos, não apenas o tamanho
de uma lista de parâmetros. Códigos universais para inteiros e sua relação
com comprimentos autodelimitados são clássicos
\citep{elias1975,coverthomas2006}. Para $n\ge1$, seja
\[
 \gamma(n)=0^{\lfloor\log_2n\rfloor}\operatorname{bin}(n)
\]
o código gama de Elias, e ponha
\[
 \zeta(z)=\begin{cases}2z,&z\ge0,\\-2z-1,&z<0.
 \end{cases}
\]
Para $r=p/q$ reduzido, $q>0$, e para o diádico escrito canonicamente como
$t=u2^{-k}$, com $u$ ímpar quando $k>0$, defina
\begin{align}
 C_{\mathbb Z}(z)&=\gamma(\zeta(z)+1),\\
 C_{\rm rat}(p/q)&=C_{\mathbb Z}(p)\,\gamma(q),
 &L_{\rm rat}(p/q)&=|C_{\rm rat}(p/q)|,\label{eq:rational-code}\\
 C_{\rm dyad}(u,k)&=C_{\mathbb Z}(u)\,\gamma(k+1),
 &L_{\rm dyad}(u,k)&=|C_{\rm dyad}(u,k)|.
 \label{eq:dyadic-code}
\end{align}
Cada componente é prefix-free e a quantidade de componentes é determinada
pelo cabeçalho do objeto; logo as concatenações abaixo são decodificáveis sem
separadores. Em particular, isso declara exatamente o que significa escrever
um diádico ``estruturalmente'' ou como racional reduzido.

\begin{lemma}[cálculo elementar de alturas]
\label{lem:height-arithmetic}
Para quaisquer $r,s\in\mathbb Q$,
\begin{equation}
 \height(r\pm s)\le2\height(r)\height(s),
 \qquad
 \height(rs)\le\height(r)\height(s).
 \label{eq:height-arithmetic}
\end{equation}
Em particular, se ambos têm altura no máximo $H$, as cotas respectivas são
$2H^2$ e $H^2$.
Para uma montagem de altura atômica máxima $H$, todo nó interno
$t_j=u_j2^{-k_j}$ satisfaz $|u_j|\le H$ e $k_j\le H$ em sua forma reduzida;
além disso,
\begin{align}
 \height(s_j-s_{j-1})&\le2H^2,\label{eq:slope-jump-height}\\
 \height(s_jt_j)&\le H2^H,\label{eq:slope-knot-height}\\
 \height\bigl(g(t_0)-s_0t_0\bigr)&\le2^{H+1}H^2.
 \label{eq:intercept-height}
\end{align}
\end{lemma}

\begin{proof}
Escreva $r=p/q$ e $s=a/b$ em forma reduzida e ponha
$R=\height(r)$, $S=\height(s)$. Antes de qualquer cancelamento,
$r\pm s=(pb\pm aq)/(qb)$ e $rs=pa/(qb)$; os numeradores têm módulos no
máximo $2RS$ e $RS$, respectivamente, e os denominadores são no máximo
$RS$. Reduzir não aumenta altura, o que prova
\eqref{eq:height-arithmetic}. Um nó é fronteira de uma célula
$J_{k,m}$, logo sua forma não reduzida tem expoente $k\le H-1$ e numerador
$m$ ou $m+1$, de módulo no máximo $H$; redução só diminui esses dados. As
cotas \eqref{eq:slope-jump-height} e \eqref{eq:slope-knot-height} seguem da
primeira parte e de $\height(t_j)\le2^H$. Finalmente, na primeira folha,
$g(t_0)-s_0t_0=y_0-s_0x_0$; como
$\height(x_0)\le2^H$, duas aplicações das cotas de produto e soma dão
\eqref{eq:intercept-height}.
\end{proof}

Fixe a ordem espacial das folhas. Os três fluxos executáveis são:
\begin{enumerate}[label=(\roman*),leftmargin=2.1em]
 \item $C_{\rm atom}(g)$: $\gamma(N)$, seguido, para cada folha, de
 $\gamma(k+1)C_{\mathbb Z}(m)C_{\rm rat}(y_0)C_{\rm rat}(s)$;
 \item $C_{\rm ReLU,str}(g)$: $\gamma(N)$, o primeiro átomo na forma ancorada
 $C_{\rm dyad}(x_0)C_{\rm rat}(y_0)C_{\rm rat}(s_0)$, e os pares
 $C_{\rm dyad}(t_j)C_{\rm rat}(s_j-s_{j-1})$;
 \item $C_{\rm ReLU,rat}(g)$: $\gamma(N)$ e a forma expandida
 $C_{\rm rat}(c)C_{\rm rat}(s_0)$, seguida dos pares
 $C_{\rm rat}(t_j)C_{\rm rat}(s_j-s_{j-1})$, onde
 $c=g(t_0)-s_0t_0$.
\end{enumerate}
Denote os comprimentos respectivos por $L_{\rm atom}$,
$L_{\rm ReLU,str}$ e $L_{\rm ReLU,rat}$. O primeiro fluxo é uma alternativa
estrutural direta à serialização por $\Gamma^\sharp$ do
\cref{thm:assembly-serialization}; ambos têm a mesma ordem, mas não a mesma
palavra.

\begin{theorem}[comparação rigorosa dos formatos]
\label{thm:coding-formats}
Os três fluxos acima são prefix-free e possuem decodificadores exatos. Para
toda montagem com $N$ folhas e altura atômica máxima $H\ge2$,
\begin{equation}
 L_{\rm atom}(g)=O(N\log H),\qquad
 L_{\rm ReLU,str}(g)=O(N\log H),\qquad
 L_{\rm ReLU,rat}(g)=O(NH),
 \label{eq:coding-format-orders}
\end{equation}
com constantes universais. A última ordem é atingida conjuntamente em $N$ e
$H$: para todo $H\ge2$ e todo $N$ par com $2\le N\le H$, considere as folhas
$J_{H-1,j}$, $0\le j<N$, com valores centrais zero e inclinações $(-1)^j$.
Essa montagem contínua tem altura $H$, e
\begin{equation}
 L_{\rm atom}=O(N\log H),\qquad
 L_{\rm ReLU,str}=O(N\log H),\qquad
 L_{\rm ReLU,rat}=\Omega(NH).
 \label{eq:coding-gap-family}
\end{equation}
Assim, a equivalência funcional montagem--ReLU preserva a ordem de bits sob o
formato estrutural fixado, mas não sob a expansão racional ingênua dos
diádicos.
\end{theorem}

\begin{proof}
O código gama é prefix-free; o mapa $\zeta$ é bijetivo entre $\mathbb Z$ e
$\Nzero$, e frações reduzidas e pares diádicos canônicos são únicos. Lendo o
cabeçalho $N$, cada decodificador sabe quantos blocos seguintes consumir. No
fluxo atômico, os blocos recuperam exatamente $(k,m,y_0,s)$ de cada folha. No
fluxo ReLU estrutural, o primeiro bloco recupera $(x_0,y_0,s_0)$ e, portanto,
$c=y_0-s_0x_0$; os demais recuperam os nós e saltos de inclinação da
\cref{eq:assembly-relu}; o terceiro fluxo já transmite esses parâmetros
expandidos. Logo os decodificadores são exatos.

Se $|z|\le A$, então $|C_{\mathbb Z}(z)|=O(\log(1+A))$, e se
$\height(r)\le A$, então $L_{\rm rat}(r)=O(\log(1+A))$. As coordenadas dos
átomos têm tamanho $O(H)$; os pares diádicos estruturais têm numerador e
expoente $O(H)$; e os saltos têm altura $O(H^2)$ pelo
\cref{lem:height-arithmetic}. Isso prova as duas primeiras cotas de
\eqref{eq:coding-format-orders}. Na forma ingênua, cada nó pode ter
denominador $2^H$, e o intercepto pode ter altura $2^{H+1}H^2$; o
\cref{lem:height-arithmetic} dá $O(H)$ bits por bloco e prova a terceira cota.

Na família final, as duas peças adjacentes ao nó
$t_j=j2^{-(H-1)}$ assumem ali o mesmo valor
$(-1)^{j-1}2^{-H}$, logo a montagem é contínua. Como $j<N\le H$, sua altura
atômica máxima é exatamente $H$. Há ao menos $(N-2)/2$ índices internos
ímpares $j$; para eles, $t_j$ já está reduzido e tem denominador
$2^{H-1}$. Cada bloco racional correspondente custa $\Omega(H)$ bits por
\eqref{eq:rational-code}, e os saltos de inclinação são não nulos. Portanto
$L_{\rm ReLU,rat}=\Omega(NH)$; os blocos atômicos e os pares diádicos
estruturais custam $O(\log H)$ cada, dando as outras duas cotas. Para $N=2$,
o intercepto $c=-2^{-H}$ fornece diretamente a mesma cota $\Omega(H)$.
Isso prova \eqref{eq:coding-gap-family}.
\end{proof}

O escopo do teorema é a comparação das palavras sob as convenções declaradas.
Complexidade de Kolmogorov, armazenamento de um codebook compartilhado e
custos de seleção ou decodificação pertencem a modelos adicionais.

\subsection{Erro conjunto sob orçamentos de folhas e altura}
\label{subsec:joint-budget}

Uma função pode admitir mais de uma representação por folhas. Para não
atribuir a ela uma altura dependente de uma escolha oculta, os dois
orçamentos serão impostos à mesma representação.

\begin{definition}[erro conjunto número--altura]
\label{def:joint-error}
Para inteiros $N,H\ge2$, seja $\Aclass_{N,H}(K)$ o conjunto das funções
$g\in\Aclass(K)$ que admitem uma representação como na
\cref{def:assembly} com no máximo $N$ folhas e com todos os átomos em
$\Library_H$. Defina
\begin{equation}
 E_{N,H}^{\Library}(f)_\infty
 :=\inf_{g\in\Aclass_{N,H}(K)}\|f-g\|_{\infty,K},
 \qquad \inf\varnothing:=+\infty.
 \label{eq:joint-error}
\end{equation}
Quando o conjunto abaixo não é vazio, ponha ainda
\begin{equation}
 H_{\varepsilon,N}^{\Library}(f)
 :=\min\{H\ge2:E_{N,H}^{\Library}(f)_\infty\le\varepsilon\}.
 \label{eq:height-for-accuracy}
\end{equation}
\end{definition}

Assim, $N$ mede a complexidade geométrica e $H$ mede, sob o calibre
\eqref{eq:atom-height}, simultaneamente resolução, posição e complexidade
aritmética. A dependência desse calibre discutida após o
\cref{lem:atom-count} continua valendo: $H$ não é uma invariante sem escolha
de exaustão.

\begin{proposition}[monotonicidade dos dois orçamentos]
\label{prop:joint-error-monotonicity}
Para todo $f\in C(K)$ e $N,H\ge2$,
\begin{equation}
 E_{N+1,H}^{\Library}(f)_\infty\le E_{N,H}^{\Library}(f)_\infty,
 \qquad
 E_{N,H+1}^{\Library}(f)_\infty\le E_{N,H}^{\Library}(f)_\infty.
 \label{eq:joint-error-monotonicity}
\end{equation}
Se
\[
 E_N^{\Library}(f)_\infty
 :=\inf\{\|f-g\|_{\infty,K}:g\in\Aclass(K)
       \text{ admite no máximo $N$ folhas}\},
\]
então $E_{N,H}^{\Library}(f)_\infty\downarrow
E_N^{\Library}(f)_\infty$ quando $H\to\infty$.
\end{proposition}

\begin{proof}
As inclusões
$\Aclass_{N,H}(K)\subseteq\Aclass_{N+1,H}(K)$ e
$\Aclass_{N,H}(K)\subseteq\Aclass_{N,H+1}(K)$ provam
\eqref{eq:joint-error-monotonicity}. Toda representação finita tem altura
máxima finita, logo a união de $\Aclass_{N,H}(K)$ sobre $H\ge2$ é exatamente
a classe usada para definir $E_N^{\Library}$. O ínfimo sobre essa união é o
limite da sequência decrescente de ínfimos.
\end{proof}

Fixe agora $R,M>0$ e considere
\begin{equation}
 \mathcal{F}_{R,M}=\{f\in C^2[0,1]:\|f\|_\infty\le R,\
                                  \ \|f''\|_\infty\le M\}.
 \label{eq:function-class}
\end{equation}

\begin{lemma}[controle uniforme da primeira derivada]
\label{lem:derivative-control}
Se $f\in\mathcal F_{R,M}$ e
\begin{equation}
 C_0:=4R+\frac M2,
 \label{eq:derivative-constant}
\end{equation}
então $\|f'\|_\infty\le C_0$.
\end{lemma}

\begin{proof}
Para $x\in[0,1]$, escolha $a\in\{0,1\}$ com $|a-x|\ge1/2$. Taylor com
resto de Lagrange dá
\begin{equation}
 |f'(x)|
 \le\frac{|f(a)|+|f(x)|}{|a-x|}+\frac{M}{2}|a-x|
 \le4R+\frac M2=C_0.
 \label{eq:derivative-control}
\end{equation}
\end{proof}

\begin{theorem}[controle simultâneo de erro, folhas, altura e bits]
\label{thm:joint-budget}
Defina
\begin{equation}
 C_H:=\max\{4(R+1),\ 2(C_0+1)\}.
 \label{eq:joint-height-constant}
\end{equation}
Para toda potência diádica $n=2^r\ge2$ e todo
$f\in\mathcal F_{R,M}$, existe uma montagem $g_n$ com exatamente $n$
folhas tal que
\begin{align}
 \|f-g_n\|_\infty
 &\le\frac{M+4}{8n^2},
 \label{eq:joint-dyadic-error}\\
 \height_{\max}(g_n)&\le C_Hn^2.
 \label{eq:joint-dyadic-height}
\end{align}
Nos três formatos do \cref{thm:coding-formats}, a mesma montagem satisfaz
\begin{equation}
 L_{\rm atom}(g_n),\ L_{\rm ReLU,str}(g_n),\
 L_{\rm ReLU,rat}(g_n)=O_{R,M}(n\log(2n)).
 \label{eq:joint-dyadic-bits}
\end{equation}
Consequentemente, para todo inteiro $N\ge2$,
\begin{equation}
 E_{N,\lceil C_HN^2\rceil}^{\Library}(f)_\infty
 <\frac{M+4}{2N^2}.
 \label{eq:joint-all-N}
\end{equation}
A construção é determinística: quando $f(t_j)$ é equidistante dos dois
pontos admissíveis da grade alternada, escolhe-se o ponto menor, ou,
equivalentemente no denominador comum de \eqref{eq:joint-common-denominator},
o menor numerador $b_j$.
\end{theorem}

\begin{proof}
Nos nós $t_j=j/n$, $0\le j\le n$, ponha $\Delta=n^{-2}$ e use as grades
alternadas
\begin{equation}
 G_j^{(n)}=\Delta\left(\mathbb Z+\frac{j\bmod2}{2}\right).
 \label{eq:joint-alternating-grid}
\end{equation}
Escolha $v_j\in G_j^{(n)}$ mais próximo de $f(t_j)$; no empate, escolha o
ponto menor. Então $|v_j-f(t_j)|\le\Delta/2$, e grades consecutivas são
disjuntas, de modo que $v_{j+1}\ne v_j$. O interpolante $g_n$ desses dados
tem, portanto, $n$ peças com inclinações racionais não nulas; cada peça é um
átomo pelo \cref{thm:dictionary}. A estimativa de interpolação e a
estabilidade convexa nos valores nodais dão
\[
 \|f-g_n\|_\infty
 \le\frac{M}{8n^2}+\frac{1}{2n^2}
 =\frac{M+4}{8n^2},
\]
o que prova \eqref{eq:joint-dyadic-error}.

Resta acompanhar a aritmética sem recorrer a uma cota genérica de soma de
racionais. Escreva, de modo comum,
\begin{equation}
 v_j=\frac{b_j}{2n^2},\qquad b_j\in\mathbb Z,
 \qquad b_j\equiv j\pmod2.
 \label{eq:joint-common-denominator}
\end{equation}
Como $|v_j|\le R+(2n^2)^{-1}$,
\begin{equation}
 |b_j|\le2Rn^2+1.
 \label{eq:joint-numerator-bound}
\end{equation}
Na folha $J_{r,j}$, o valor central e a inclinação são
\begin{equation}
 y_{0,j}=\frac{b_j+b_{j+1}}{4n^2},
 \qquad
 s_j=\frac{b_{j+1}-b_j}{2n}.
 \label{eq:joint-atom-parameters}
\end{equation}
O numerador de $s_j$ é ímpar, portanto a segunda fração já está reduzida e é
não nula. Pelo \cref{lem:derivative-control},
\begin{equation}
 |b_{j+1}-b_j|
 =2n^2|v_{j+1}-v_j|
 \le2n^2\left(\frac{C_0}{n}+\frac1{n^2}\right)
 =2C_0n+2.
 \label{eq:joint-slope-numerator}
\end{equation}
Reduzir $y_{0,j}$ só diminui sua altura; logo
\begin{equation}
 \height(y_{0,j})\le4(R+1)n^2,
 \qquad
 \height(s_j)\le2(C_0+1)n.
 \label{eq:joint-parameter-heights}
\end{equation}
Além disso, $r+1\le n$ e $|j|+1\le n$. A definição
\eqref{eq:atom-height} e \eqref{eq:joint-height-constant} provam
\eqref{eq:joint-dyadic-height}.

Para os bits, no formato ReLU expandido o intercepto é $c=v_0$,
$\height(c)\le2(R+1)n^2$; cada nó interno tem altura no máximo $n$; e, de
\eqref{eq:joint-slope-numerator},
\[
 \height(s_j-s_{j-1})\le4(C_0+1)n.
\]
Assim, tanto os parâmetros atômicos quanto os parâmetros ReLU desta
construção têm altura polinomial em $n$. Nos códigos concretos
\eqref{eq:rational-code}--\eqref{eq:dyadic-code}, cada bloco usa
$O_{R,M}(\log(2n))$ bits; há $O(n)$ blocos em qualquer dos três fluxos. Isso
prova \eqref{eq:joint-dyadic-bits}. Observe que a cota geral
$O(NH)$ do formato racional ingênuo não é saturada por esta família: seus
nós têm denominador apenas $n$, e não $2^H$.

Por fim, dado $N\ge2$, tome $n=2^{\lfloor\log_2N\rfloor}$. Então
$N/2<n\le N$. A montagem já construída usa no máximo $N$ folhas, tem altura
no máximo $C_Hn^2\le C_HN^2$ e erro menor que
$(M+4)/(2N^2)$, provando \eqref{eq:joint-all-N}.
\end{proof}

\begin{corollary}[altura suficiente para precisão prescrita]
\label{cor:height-for-epsilon}
Para $\varepsilon>0$, seja
\begin{equation}
 n_\varepsilon
 =2^{\left\lceil\log_2\max\left\{2,
        \sqrt{\frac{M+4}{8\varepsilon}}\right\}\right\rceil},
 \qquad
 H_\varepsilon=\lceil C_Hn_\varepsilon^2\rceil.
 \label{eq:epsilon-schedule}
\end{equation}
Então, uniformemente para $f\in\mathcal F_{R,M}$,
\begin{align}
 E_{n_\varepsilon,H_\varepsilon}^{\Library}(f)_\infty
 &\le\varepsilon,\label{eq:epsilon-joint-error}\\
 n_\varepsilon&=O_{R,M}(\varepsilon^{-1/2}),
 &H_\varepsilon&=O_{R,M}(\varepsilon^{-1}),
 \label{eq:epsilon-joint-height}\\
 L_{\rm atom},\ L_{\rm ReLU,str},\ L_{\rm ReLU,rat}
 &=O_{R,M}\!\left(\varepsilon^{-1/2}
                   \log\frac{2}{\varepsilon}\right).
 \label{eq:epsilon-joint-bits}
\end{align}
Em particular,
$H_{\varepsilon,n_\varepsilon}^{\Library}(f)\le H_\varepsilon$.
\end{corollary}

\begin{proof}
A escolha de $n_\varepsilon$ torna o lado direito de
\eqref{eq:joint-dyadic-error} no máximo $\varepsilon$. A menor potência de
dois acima de um número positivo é menor que duas vezes esse número, salvo o
piso fixo $2$; as três ordens seguem então do \cref{thm:joint-budget}.
\end{proof}

O expoente $H_\varepsilon=O(\varepsilon^{-1})$ é mais forte que a estimativa
preliminar $O(\varepsilon^{-3/2})$: o ganho vem do denominador diádico comum
em \eqref{eq:joint-common-denominator} e do controle de diferenças adjacentes
por $\|f'\|_\infty$. Trata-se de uma cota superior construtiva para a altura
mínima de \eqref{eq:height-for-accuracy}, não de uma cota inferior ótima.

Para tornar as quatro quantidades auditáveis, a
\cref{tab:joint-budget-pt} usa $f_\star(x)=4x(x-1)$, $R=1$ e $M=8$. Nesse
caso $C_0=8$, $C_H=18$, o erro exato da construção é
$49/(64N^2)$ e a altura efetiva é $4N^2$. Os três comprimentos são contados
bit a bit pelos codificadores distribuídos.

\begin{table}[t]
\centering
\small
\resizebox{\linewidth}{!}{\begin{tabular}{rrrrrrrr}
\toprule
$N$ & $e_{\rm cert}$ & $e_{\rm exato}$ & $\height_{\max}$ & $H_N$ & $L_{\rm atom}$ & $L_{\rm str}$ & $L_{\rm rat}$\\
\midrule
2 & $\frac{3}{8}$ & $\frac{49}{256}$ & 16 & 72 & 77 & 59 & 37\\
4 & $\frac{3}{32}$ & $\frac{49}{1024}$ & 64 & 288 & 195 & 115 & 89\\
8 & $\frac{3}{128}$ & $\frac{49}{4096}$ & 256 & 1152 & 513 & 235 & 211\\
16 & $\frac{3}{512}$ & $\frac{49}{16384}$ & 1024 & 4608 & 1243 & 495 & 497\\
32 & $\frac{3}{2048}$ & $\frac{49}{65536}$ & 4096 & 18432 & 2921 & 1067 & 1159\\
64 & $\frac{3}{8192}$ & $\frac{49}{262144}$ & 16384 & 73728 & 6731 & 2329 & 2669\\
\bottomrule
\end{tabular}
}
\caption{Construção do \cref{thm:joint-budget} para
$f_\star(x)=4x(x-1)$. $e_{\rm cert}$ é a cota
\eqref{eq:joint-dyadic-error}, $e_{\rm exato}$ é obtido em aritmética
racional nos extremos e no ponto crítico de cada folha, e $H_N=18N^2$ é a
cota provada.}
\label{tab:joint-budget-pt}
\end{table}

\subsection{Entropia de uma classe limitada de \texorpdfstring{$C^2$}{C2}}

Se $N_\varepsilon(\mathcal{F}_{R,M})$ é o menor número de bolas de raio $\varepsilon$ em norma uniforme que cobre a classe, escreva
\[
 \mathcal{H}_\varepsilon(\mathcal{F}_{R,M})=\log_2N_\varepsilon(\mathcal{F}_{R,M}).
\]
Essa é a entropia métrica de Kolmogorov \citep{kolmogorov1959,kolmogorov1961,lorentz1996}; a ordem $\varepsilon^{-1/2}$ para classes com duas derivadas limitadas é clássica, e a demonstração abaixo é incluída apenas para tornar explícitas as constantes usadas no codebook.

\begin{theorem}[entropia e codebook explícito]\label{thm:entropy}
Existem $\varepsilon_0,c,C,C_1>0$, dependentes apenas de $R$ e $M$, tais que, para $0<\varepsilon<\varepsilon_0$,
\begin{equation}
 c\varepsilon^{-1/2}
 \le\mathcal{H}_\varepsilon(\mathcal{F}_{R,M})
 \le C\varepsilon^{-1/2}.
 \label{eq:entropy-order}
\end{equation}
Existe também um codebook finito $\mathcal{C}_\varepsilon\subset\Aclass([0,1])$, composto somente por montagens de átomos de $\Library$, com distorção no máximo $\varepsilon$. Quando o codebook e seu decodificador são informação compartilhada, o comprimento de seu índice global de comprimento fixo é
\begin{equation}
 L_\varepsilon^{\mathrm{global}}(\mathcal C_\varepsilon)
 :=\left\lceil\log_2|\mathcal{C}_\varepsilon|\right\rceil
 \le C_1\varepsilon^{-1/2}.
 \label{eq:code-length}
\end{equation}
Logo esse codebook atinge em ordem a entropia métrica da classe em cardinalidade e em bits de índice global.
\end{theorem}

\begin{proof}
\emph{Cota superior em \eqref{eq:entropy-order}.} Tome
\[
 N=\left\lceil\sqrt{M/\varepsilon}\right\rceil+1,\qquad h=1/N,
\]
e seja $P_Nf$ o interpolante nos nós $j/N$. A estimativa local \eqref{eq:interpolation-remainder}, válida para intervalos arbitrários, fornece $\|f-P_Nf\|_\infty\le Mh^2/8\le\varepsilon/8$. Quantize cada $v_j=f(j/N)$ no múltiplo mais próximo de $\Delta=\varepsilon/4$, obtendo $v'_j$ com $|v'_j-v_j|\le\varepsilon/8$. O interpolante quantizado $g$ satisfaz $\|f-g\|_\infty\le\varepsilon/4<\varepsilon$.

Resta contar os vetores quantizados que podem ocorrer. Há $O(\varepsilon^{-1})$ escolhas para $v'_0$. Dado $v'_0$, \eqref{eq:derivative-control} implica
\[
 |v'_1-v'_0|\le C_0h+\varepsilon/4,
\]
e portanto há $O((N\varepsilon)^{-1})=O(\varepsilon^{-1/2})$ escolhas para $v'_1$. Para $1\le j\le N-1$, a identidade integral
\[
 v_{j+1}-2v_j+v_{j-1}
 =\int_0^h\int_0^h
 f''((j-1)h+s+t)\dd s\dd t
\]
dá módulo no máximo $Mh^2\le\varepsilon$. A quantização altera a diferença segunda em no máximo $4(\varepsilon/8)=\varepsilon/2$. Dados $v'_{j-1}$ e $v'_j$, o próximo valor está, pois, em um intervalo de comprimento $3\varepsilon$, que contém no máximo uma constante universal $K$ de pontos da grade de passo $\varepsilon/4$. Consequentemente, o número total de vetores é
\[
 O(\varepsilon^{-1})\,O(\varepsilon^{-1/2})\,K^N
 =\exp\bigl(O(\varepsilon^{-1/2})\bigr),
\]
e a cota superior segue ao tomar $\log_2$.

\emph{Cota inferior em \eqref{eq:entropy-order}.} Fixe $\psi\in C_c^2(0,1)$ com $\|\psi\|_\infty=1$. Ponha $A=4\varepsilon$ e
\[
 w=\sqrt{\frac{A\|\psi''\|_\infty}{M}}.
\]
Para $\varepsilon$ pequeno, $A\le R$ e cabem $m=\lfloor1/w\rfloor\asymp\varepsilon^{-1/2}$ intervalos consecutivos de largura $w$ em $[0,1]$. Se $a_j=(j-1)w$, defina
\[
 \varphi_j(x)=A\psi\!\left(\frac{x-a_j}{w}\right),\qquad
 f_S=\sum_{j\in S}\varphi_j,\quad S\subset\{1,\ldots,m\}.
\]
Os suportes são disjuntos, $\|f_S\|_\infty\le A\le R$ e $\|f_S''\|_\infty\le M$. Se $S\neq T$, então $\|f_S-f_T\|_\infty=A=4\varepsilon$. Logo os $2^m$ elementos formam um empacotamento $4\varepsilon$-separado; uma bola de raio $\varepsilon$ contém no máximo um deles. Assim, $N_\varepsilon\ge2^m$ e $\mathcal{H}_\varepsilon\ge m\ge c\varepsilon^{-1/2}$.

\emph{Codebook de montagens.} Seja $N=2^k$ a menor potência de dois com $M/(8N^2)\le\varepsilon/2$. Pela minimalidade da potência de dois, $N=\Theta(\varepsilon^{-1/2})$ e $h=N^{-1}=\Theta(\varepsilon^{1/2})$. Defina
\[
 \delta=2^{-\lceil\log_2(1/\varepsilon)\rceil},
 \qquad \varepsilon/2<\delta\le\varepsilon,
\]
e, no nó $t_j=j/N$, use a grade finita alternada
\begin{equation}
 G_j=\delta\left(\mathbb{Z}+\frac{j\bmod2}{2}\right)
 \cap[-R-\delta,R+\delta].
 \label{eq:alternating-grid}
\end{equation}
Para cada $f$, escolha em $G_j$ um ponto $v_j$ mais próximo de $f(t_j)$, com desempate determinístico. Então $|v_j-f(t_j)|\le\delta/2$. As grades de paridades opostas são disjuntas, de modo que $v_{j+1}\neq v_j$. Como $N$ e $\delta$ são diádicos, o interpolante $g$ dos dados $(t_j,v_j)$ é uma montagem de átomos. Além disso,
\[
 \|f-g\|_\infty
 \le\|f-P_kf\|_\infty+\|P_kf-g\|_\infty
 \le\frac{M}{8N^2}+\frac\delta2\le\varepsilon.
\]
Defina $\mathcal C_\varepsilon$ como a imagem desse mapa de quantização.
É exatamente o mesmo conjunto finito de interpolantes; resta apenas contá-lo
usando a restrição de curvatura, em vez de contar todos os vetores da grade.
Há $O(\varepsilon^{-1})$ escolhas para $v_0$. Dado $v_0$, o controle
\eqref{eq:derivative-control} e os erros de quantização dão
\[
 |v_1-v_0|\le C_0h+\delta,
\]
de modo que há $O(h/\delta+1)=O(\varepsilon^{-1/2})$ escolhas para $v_1$.
Para $1\le j\le N-1$, a identidade de diferença segunda já usada acima e
$Mh^2\le4\varepsilon$ implicam
\[
 |v_{j+1}-2v_j+v_{j-1}|
 \le Mh^2+2\delta\le6\varepsilon.
\]
Fixados $v_{j-1}$ e $v_j$, o valor $v_{j+1}$ fica, portanto, confinado a um
intervalo de comprimento $12\varepsilon$. Os pontos de $G_{j+1}$ têm passo
$\delta>\varepsilon/2$, logo esse intervalo contém no máximo uma constante
universal $K$ de candidatos; por exemplo, $K=25$ serve. Consequentemente,
\[
 |\mathcal C_\varepsilon|
 \le C\varepsilon^{-1}\,C'\varepsilon^{-1/2}\,K^{N-1},
\]
e, como $N=O(\varepsilon^{-1/2})$,
\[
 \log_2|\mathcal C_\varepsilon|
 \le O(\log(1/\varepsilon))+N\log_2K
 =O(\varepsilon^{-1/2}).
\]
Isso prova \eqref{eq:code-length} sem alterar os átomos, as grades alternadas
ou a garantia de distorção.
\end{proof}

\subsection{Cinco recursos distintos: cobertura, bits, serialização, seleção e decodificação}

Para impedir que uma afirmação informacional seja interpretada como uma
afirmação algorítmica, fixemos as noções usadas daqui em diante. Seja
$(X,d)$ um espaço métrico, $\mathcal F\subset X$ totalmente limitado e
$\varepsilon>0$.

\begin{definition}[recursos de um codec]\label{def:coding-resources}
Um $\varepsilon$-codec para $\mathcal F$ é uma tripla
$(\mathcal W,\mathsf{Enc},\mathsf{Dec})$, na qual $\mathcal W$ é um conjunto
finito de palavras binárias, $\mathsf{Enc}:\mathcal F\to\mathcal W$ e
$\mathsf{Dec}:\mathcal W\to X$ satisfazem
\[
 \sup_{f\in\mathcal F}d\bigl(f,\mathsf{Dec}(\mathsf{Enc}(f))\bigr)
 \le\varepsilon.
\]
O decodificador é compartilhado e seu armazenamento não é contado no
comprimento das palavras. Distinguem-se:
\begin{align*}
 \mathcal H_\varepsilon(\mathcal F)
  &=\log_2N_\varepsilon(\mathcal F),
  &&\text{entropia métrica;}\\
 L_\varepsilon^{\mathrm{dec}}(\mathcal F)
  &=\inf_{\text{$\varepsilon$-codecs prefix-free}}
    \max_{w\in\mathcal W}|w|,
  &&\text{menor comprimento máximo decodificável;}\\
 L_\varepsilon^{\mathrm{global}}(\mathcal C_\varepsilon)
  &=\left\lceil\log_2|\mathcal C_\varepsilon|\right\rceil,
  &&\text{índice fixo em um codebook compartilhado.}
\end{align*}
Para um codebook de montagens e uma convenção de serialização fixada, ponha
\begin{align*}
 L^{\mathrm{folhas}}(g\mid K)
   &=L_{\mathrm{arv}}(g\mid K)
     +\sum_{J\in\Lambda(g)}L_{\mathrm{addr}}(g|_J),\\
 L_\varepsilon^{\mathrm{folhas}}(\mathcal C_\varepsilon\mid K)
   &=\max_{g\in\mathcal C_\varepsilon}L^{\mathrm{folhas}}(g\mid K).
\end{align*}
Por fim, depois de fixados um modelo de computação $\mathfrak M$, uma
representação de entrada para $f$ e o codec, $T_\varepsilon^{\mathrm{enc}}$
é o custo de selecionar/produzir a palavra e $T_\varepsilon^{\mathrm{dec}}$ é
o custo de reconstruí-la. Sem essas três escolhas, os dois custos temporais
não são quantidades matematicamente definidas.
\end{definition}

O símbolo ``global'' significa comprimento do \emph{índice} em uma tabela
compartilhada, não o número de bits necessário para armazenar a tabela inteira.
O símbolo ``dec'' em $L_\varepsilon^{\mathrm{dec}}$ qualifica a palavra como
decodificável; ele não representa tempo de decodificação.

\begin{proposition}[limite informacional para representação decodificável]
\label{prop:description-lower}
Todo $\varepsilon$-codec de comprimento fixo $\ell$ para $\mathcal F$ satisfaz
\begin{equation}
 \ell\ge\left\lceil\mathcal H_\varepsilon(\mathcal F)\right\rceil.
 \label{eq:description-lower}
\end{equation}
A mesma desigualdade vale se $\mathcal W$ é prefix-free e $\ell$ é o maior
comprimento de palavra. Em particular,
\[
 L_\varepsilon^{\mathrm{dec}}(\mathcal F)
 \ge\left\lceil\mathcal H_\varepsilon(\mathcal F)\right\rceil.
\]
Sob convenções que acrescentem ou removam apenas um cabeçalho fixo, a forma
invariante é $\ell\ge\mathcal H_\varepsilon(\mathcal F)-O(1)$.
\end{proposition}

\begin{proof}
As reconstruções $\mathsf{Dec}(\mathcal W)$ formam uma cobertura de
$\mathcal F$ por bolas de raio $\varepsilon$. Portanto,
\[
 N_\varepsilon(\mathcal F)
 \le |\mathsf{Dec}(\mathcal W)|\le|\mathcal W|.
\]
No caso de comprimento fixo, $|\mathcal W|\le2^\ell$. No caso prefix-free, a
desigualdade de Kraft e $|w|\le\ell$ dão
\[
 |\mathcal W|2^{-\ell}
 \le\sum_{w\in\mathcal W}2^{-|w|}\le1,
\]
e novamente $|\mathcal W|\le2^\ell$. Tomar logaritmos prova
\eqref{eq:description-lower}. Essa separação entre entropia e complexidade de
descrição é a que se usa em limites informacionais para esquemas quantizados
\citep{bolcskei2019}.
\end{proof}

Aplicada ao \cref{thm:entropy}, a proposição produz a cadeia precisa
\begin{equation}
 \mathcal H_\varepsilon(\mathcal F_{R,M})
 \le L_\varepsilon^{\mathrm{dec}}(\mathcal F_{R,M})
 \le L_\varepsilon^{\mathrm{global}}(\mathcal C_\varepsilon)
 =O(\varepsilon^{-1/2}).
 \label{eq:complexity-chain}
\end{equation}
Como a entropia também é $\Omega(\varepsilon^{-1/2})$, o codec é ótimo em
ordem quanto à cardinalidade, ao comprimento decodificável e aos bits do
índice global. Os custos $T_\varepsilon^{\mathrm{enc}}$ e
$T_\varepsilon^{\mathrm{dec}}$ constituem recursos computacionais separados.

Já a serialização operacional folha a folha obedece ao
\cref{thm:assembly-serialization}. Na construção do \cref{thm:entropy}, as
alturas dos parâmetros são polinomiais em $1/\varepsilon$, e por isso pode
reaparecer um custo $O(\log(1/\varepsilon))$ por folha. Esse é um custo do
formato de serialização, não da cardinalidade necessária para cobrir
$\mathcal F_{R,M}$.

\subsection{Catálogos finitos, minimizadores e truncagem}

Densidade não garante a existência de minimizadores sobre uma biblioteca
infinita. O fenômeno já ocorre para uma única avaliação.

\begin{proposition}\label{prop:no-minimizer}
Para $x=1/4$ e $y=\sqrt2$,
\[
 \inf\{|y-\ell_n(x)|:n\in\Nzero,\ x\in I_n^x\}=0,
\]
mas nenhum átomo realiza o ínfimo.
\end{proposition}

\begin{proof}
Para todo átomo, $a_n=y_{0,n}-s_nx_{0,n}\in\mathbb{Q}$ e
$b_n=s_n\in\mathbb{Q}$; como $x=1/4$ é racional,
$\ell_n(x)\in\mathbb{Q}$. Portanto, $\ell_n(x)\neq\sqrt2$.

Por outro lado, dado $\eta>0$, escolha $q\in\mathbb{Q}$ com
$|q-\sqrt2|<\eta$. Use a célula $J_{2,1}=[1/4,1/2]$, de centro $x_0=3/8$,
escolha $s=1$ e $y_0=q-s(x-x_0)=q+1/8\in\mathbb{Q}$. O
\cref{thm:dictionary} fornece um átomo nessa célula com
$\ell_n(1/4)=q$. Assim, os erros tendem a zero, mas nenhum é zero.
\end{proof}

A correção matemática é truncar simultaneamente geometria e aritmética. Fixe
uma raiz diádica $K=J_{k_0,m_0}$ e seja $\mathfrak T_D(K)$ o conjunto das
partições obtidas por árvores binárias diádicas completas de profundidade
relativa no máximo $D$. Uma \emph{montagem representada} é um par
$R=(\Lambda,n)$ formado por $\Lambda\in\mathfrak T_D(K)$ e um endereço
$n(J)$ para cada folha, com suporte exatamente $J$ e concordância dos valores
nas fronteiras.

\begin{definition}[catálogo truncado]\label{def:finite-catalog}
Para uma montagem representada, ponha
\[
 \height_{\max}(R)=\max_{J\in\Lambda}
 \max\{2,k(J)+1,|m(J)|+1,\height(y_{0,J}),\height(s_J)\}.
\]
Se
$H\ge H_K:=\max\{2,k_0+1,|m_0|+1\}$, defina
\begin{align}
 \mathfrak R_{D,H}^{\Library}(K)
   &=\{R=(\Lambda,n):\Lambda\in\mathfrak T_D(K),\
       \height_{\max}(R)\le H\},\label{eq:represented-catalog}\\
 \Sigma_{D,H}^{\Library}(K)
   &=\{\operatorname{Eval}(R):R\in
       \mathfrak R_{D,H}^{\Library}(K)\}\subset\Aclass(K).
 \label{eq:finite-catalog}
\end{align}
Representações distintas que avaliam na mesma função correspondem, portanto,
a um único elemento de $\Sigma_{D,H}^{\Library}(K)$.
\end{definition}

Para fixar o desempate, serialize a árvore em pré-ordem com $0$ para nó
interno e $1$ para folha e, depois, escreva os endereços
$n^\sharp(J)+1$ das folhas em ordem espacial com o código gama de Elias. O
esqueleto informa quantas palavras devem ser lidas, logo o fluxo
$\operatorname{ser}_\sharp(R)$ é autodelimitado. Ordene os fluxos por
\emph{shortlex}: primeiro pelo comprimento e depois lexicograficamente. Para
cada função, retenha o menor fluxo entre todas as suas representações no
catálogo.

\begin{theorem}[finitude, existência e desempate canônico]
\label{thm:finite-catalog-minimizer}
O catálogo $\Sigma_{D,H}^{\Library}(K)$ é finito e não vazio. Mais
precisamente, se
\[
 T_0=1,\qquad T_{d+1}=1+T_d^2,
\]
então
\begin{equation}
 1\le |\Sigma_{D,H}^{\Library}(K)|
 \le |\mathfrak R_{D,H}^{\Library}(K)|
 \le T_D|\Library_H|^{2^D}<\infty.
 \label{eq:finite-catalog-bound}
\end{equation}
Consequentemente, para qualquer funcional
$\mathcal J:\Sigma_{D,H}^{\Library}(K)\to\mathbb R$, existe
\begin{equation}
 g^\star\in\argmin_{g\in\Sigma_{D,H}^{\Library}(K)}\mathcal J(g).
 \label{eq:finite-argmin}
\end{equation}
Entre todos os minimizadores, a regra do menor fluxo shortlex seleciona uma
única função e uma única representação canônica.
\end{theorem}

\begin{proof}
A recorrência de $T_D$ conta todas as formas ordenadas de árvore: ou a raiz é
folha, ou seus dois filhos recebem independentemente uma forma de profundidade
no máximo $D-1$. Cada árvore tem no máximo $2^D$ folhas e cada folha recebe
um rótulo de $\Library_H$. Ignorar as restrições de suporte e continuidade só
aumenta a contagem, o que prova a cota superior. A árvore reduzida à raiz pode
ser rotulada, por exemplo, com $y_0=0$ e $s=1$, logo o catálogo não é vazio.
A imagem de um conjunto finito por $\operatorname{Eval}$ é finita; qualquer
função real em um conjunto finito não vazio atinge seu mínimo.

Resta o desempate. Há um número finito de representações de cada função, e a
ordem shortlex é total; portanto existe um menor fluxo. A serialização é
injetiva sobre representações, pois recupera a árvore e cada endereço de
folha. Assim, dois candidatos distintos não compartilham o mesmo fluxo, e o
menor minimizador é único.
\end{proof}

\begin{proposition}[consistência sob truncagem]
\label{prop:catalog-consistency}
Se $D\le D'$ e $H\le H'$, então
\begin{equation}
 \Sigma_{D,H}^{\Library}(K)
 \subseteq\Sigma_{D',H'}^{\Library}(K).
 \label{eq:catalog-nesting}
\end{equation}
Para um funcional fixado na união dos catálogos, os valores mínimos são, por
isso, não crescentes em $D$ e $H$. Se o funcional possui um minimizador global
nessa união, o minimizador global de menor fluxo shortlex e sua representação
canônica pertencem a algum catálogo finito; a partir desse par de truncagens,
o seletor canônico estabiliza.
\end{proposition}

\begin{proof}
Uma representação admissível sob $(D,H)$ continua admissível quando os dois
limites aumentam, o que prova \eqref{eq:catalog-nesting} e a monotonicidade dos
mínimos. A ordem shortlex bem ordena as palavras binárias finitas. Logo, entre
todas as representações de todos os minimizadores globais, existe uma menor;
ela tem profundidade e altura finitas. Qualquer catálogo que já a contenha não
pode ser superado posteriormente nem no valor do funcional nem no desempate.
\end{proof}

\begin{algorithm}[enumeração exata do catálogo]\label{alg:finite-enumeration}
Para parâmetros pequenos, gere recursivamente todas as formas de
$\mathfrak T_D(K)$; em cada folha, liste somente os átomos de $\Library_H$ com
o suporte exigido; ao unir duas subárvores, descarte pares cujos valores na
fronteira diferem. Reduza representações funcionalmente idênticas ao menor
fluxo $\operatorname{ser}_\sharp$, avalie $\mathcal J$ em aritmética racional
exata e aplique o desempate shortlex.
\end{algorithm}

Esse algoritmo exaustivo realiza operacionalmente o
\cref{thm:finite-catalog-minimizer}. A cota
\eqref{eq:finite-catalog-bound} delimita seu uso a instâncias pequenas;
eficiência assintótica, estabilidade em relação aos dados e comportamento
perto de empates são recursos adicionais \citep{christof2021}.
